% Image credits for the photographs and AI illustrations of Book 3 (University, Year 1).
% Machine-readable license records live in images/book3/CREDITS.md; keep the
% two files in sync. The Book 3 writer owns this file: add one \item per
% photograph (inside the itemize, by chapter) as photographs land.
\chapter*{Bronvermelding van afbeeldingen}

De tekeningen, kaarten en grafieken in dit boek zijn origineel; kaarten en
grafieken die op gepubliceerde gegevens zijn gebaseerd, vermelden hun bron in
het bijschrift of in deze lijst. Foto's worden gebruikt onder hun respectieve
vrije licenties, met dank aan hun makers; volledige bronlinks en
licentie-URL's voor elke foto staan in \texttt{images/book3/CREDITS.md} in de
repository van het boek.

\bigskip
Naast de foto's en de originele tekeningen zijn sommige figuren door AI
gegenereerde illustraties, gemaakt met OpenAI-beeldgeneratie op basis van
prompts die door de redacteuren van het boek zijn geschreven, en vóór opname
gecontroleerd op wetenschappelijke juistheid. Ze tonen uitsluitend generieke
scènes, nooit een echte plaats, een echt persoon of een echt exemplaar. De
volledige prompt voor elke gegenereerde afbeelding staat in
\texttt{images/book3/ai/PROMPTS.md} in de repository.

\bigskip
\noindent Foto's en gegevens, per hoofdstuk:
\begin{itemize}\small
  \item Hfst.~1 (sferische astronomie): sterrensporen boven het Paranal-observatorium, R.~Wesson/ESO, CC BY 4.0.
  \item Hfst.~2 (telescopen, detectoren en fotometrie): zes segmenten van de primaire spiegel van de James Webb-ruimtetelescoop vóór de cryogene tests, NASA/MSFC/David Higginbotham, publiek domein.
  \item Hfst.~3 (vergelijkende planetologie): de aardachtige planeten op schaal, NASA; Venus met de radar van Magellan, NASA/JPL; Mars van de Viking~1-baan, NASA/USGS; Europa, NASA/JPL/DLR; Io, NASA/JPL/University of Arizona; allemaal publiek domein.
  \item Hfst.~4 (geologische kaarten): William Smith, \emph{A Delineation of the Strata of England and Wales, with part of Scotland}, 1815; publiek domein (Wikimedia Commons).
  \item Hfst.~5 (geologische doorsneden): Hutton's ongelijkvormigheid bij Siccar Point, PawełS, CC BY-SA 4.0 (Wikimedia Commons).
  \item Hfst.~6 (veldwerk- en laboratoriummethoden): een biotietgraniet in dunne slijp tussen gekruiste polarisatoren, Darla Sondrol, CC0 (Wikimedia Commons).
  \item Hfst.~7 (gesteentevormende mineralen): monsters van olivijn, kwarts, biotiet, almandien en hoornblende, Robert M.~Lavinsky, CC BY-SA 3.0; labradoriet, Twyla Baker, CC0; calcietspaltrhomboëders, Giovanna Canu en Eva Santini, CC BY 4.0 (allemaal Wikimedia Commons).
  \item Hfst.~8 (optische mineralogie): een gabbro in dunne slijp met één polarisator en tussen gekruiste polarisatoren, Chmee2, CC BY-SA 3.0; Auguste Michel-Lévy als student, publiek domein (Wikimedia Commons). De interferentiekleuren zijn berekend uit de kleurvergelijkingsfuncties van de CIE 1931 en de standaardlichtbron D65 van de CIE, \copyright{} CIE, CC BY-SA 4.0.
  \item Hfst.~9 (de gelaagde aarde): Andrija Mohorovičić, publiek domein; Inge Lehmann in 1932, foto van Even Neuhaus, publiek domein (Wikimedia Commons). PREM: A.~M.~Dziewonski en D.~L.~Anderson (1981), zoals opgenomen in ObsPy.
  \item Hfst.~10 (de warmtebalans van de aarde): wereldwijde warmtestroom van J.~H.~Davies en D.~R.~Davies (2010), \emph{Solid Earth}, CC BY 3.0; manteladiabaat van T.~Katsura (2022), \emph{JGR Solid Earth}, CC BY.
  \item Hfst.~11 (isostasie en zwaartekrachtsanomalieën): het slingerapparaat van Vening Meinesz, Broodt, CC BY-SA 4.0 (Wikimedia Commons). Hoogten van ETOPO1 (NOAA National Centers for Environmental Information), om de graad bemonsterd.
  \item Hfst.~12 (plaatbeweging op een bol): een permanent station voor satellietpositionering, Alpsdake, CC BY-SA 4.0 (Wikimedia Commons). Plaatgrenzen van PB2002 (P.~Bird, 2003; GeoJSON door H.~Ahlenius, ODC-By 1.0); kustlijnen van Natural Earth (publiek domein).
  \item Hfst.~13 (slenken, ruggen en passieve randen): open scheuren bij Þingvellir, IJsland, Olga Ernst, CC BY-SA 4.0; kussenlava's op de Galápagos-rift, NOAA Okeanos Explorer Program, CC BY 2.0 (beide Wikimedia Commons). Dieptes van ETOPO1 (NOAA National Centers for Environmental Information) en ouderdommen van de oceaanbodem van M.~Seton en anderen (2020), EarthByte, CC BY 4.0.
  \item Hfst.~14 (actieve randen): Sarychev Peak in eruptie, gezien vanuit het Internationaal Ruimtestation, NASA, publiek domein (Wikimedia Commons). Aardbevingen uit de USGS ANSS Comprehensive Catalog (publiek domein); vulkanen van het Smithsonian Institution Global Volcanism Program; plaatgrenzen van PB2002 (P.~Bird, 2003; GeoJSON door H.~Ahlenius, ODC-By 1.0); kusten van Natural Earth (publiek domein); ouderdommen van de oceaanbodem van M.~Seton en anderen (2020), EarthByte, CC BY 4.0.
  \item Hfst.~15 (hotspots en intraplaatvulkanisme): een lavafontein bij Pu\textquoteleft u \textquoteleft O\textquoteleft o, Kilauea, G.~E.~Ulrich, U.S. Geological Survey, publiek domein; de West-Ghats bij Matheran, Nicholas (Nichalp), CC BY-SA 2.5 (beide Wikimedia Commons). Posities, spooraziumuts en snelheden van hotspots uit tabel~1 van het elektronische supplement van W.~J.~Morgan en J.~P.~Morgan (2007); dieptes van ETOPO1 (NOAA National Centers for Environmental Information); plaatgrenzen van PB2002 (P.~Bird, 2003; GeoJSON door H.~Ahlenius, ODC-By 1.0); kusten van Natural Earth (publiek domein).
  \item Hfst.~16 (spanning, rek en reologie): een hogedrukcel voor triaxiale proeven om de sterkte van gesteente te meten, Darryl Peroni, CSIRO, CC BY 3.0 (Wikimedia Commons). Kruipparameters van nat kwartsiet (Gleason en Tullis, 1995) en droge olivijn (Hirth en Kohlstedt, 2003), zoals samengesteld door M.~A.~Lopez-Sanchez in \emph{strength\_envelopes} (MPL-2.0).
  \item Hfst.~17 (breuken, plooien en structuren): een slickenside bij Corona Heights, San Francisco, Dextersinister, CC BY-SA 4.0; breccie van de Keystone-overschuiving, Nevada, Jstuby, publiek domein; chevronplooien bij Millook, Cornwall, Smalljim, CC BY 3.0; een antiklinaal met axiaalvlakse splijting nabij Potes, Cantabria, Mikenorton, CC BY-SA 3.0 (allemaal Wikimedia Commons).
  \item Hfst.~18 (de aardbevingsbron): land en hekken die door de breuk van 1906 in Marin County zijn gescheiden, \emph{Popular Science Monthly} vol.~69 (1906), publiek domein; Wallace Creek verschoven over de San Andreas-breuk, R.~E.~Wallace, U.S. Geological Survey, publiek domein (beide Wikimedia Commons). Aardbevingen uit de USGS ANSS Comprehensive Catalog (publiek domein); het focale mechanisme van de aardbeving van Tohoku in 2011 uit de W-fase-momententensor van de USGS (publiek domein).
  \item Hfst.~19 (verwering en kleien): verweerd graniet en granietzand bij Targassonne, Pyreneeën, Pinpin, CC~BY~2.5; bauxiet op kaolinitisch zandsteen bij Weipa, Australië, Werner Schellmann, CC~BY-SA~2.5; kaolinietkristallen onder de elektronenmicroscoop, U.S. Geological Survey, publiek domein (allemaal Wikimedia Commons). Ionstralen van R.~D.~Shannon (1976).
  \item Hfst.~20 (erosie en transport): de overstromingspluim van de Mississippi van het MODIS-instrument op Terra, 4~maart 2018, NASA Earth Observatory-beeld van Joshua Stevens, publiek domein (Wikimedia Commons).
  \item Hfst.~21 (klastische sedimenten en hun milieus): kruisgelaagdheid in de Aztec-zandsteen, Nevada, Chip Carroon, Bureau of Land Management, CC~BY-SA~2.0; de delta van de Lena vanuit Landsat~7, U.S. Geological Survey, publiek domein; een turbidiet met de onderdelen van de Bouma-reeks, Becke-Oese, Jo Weber, CC~BY-SA~3.0 (allemaal Wikimedia Commons).
  \item Hfst.~22 (carbonaat-, kiezel- en verdampingssedimenten): Bora Bora vanuit het Internationaal Ruimtestation, NASA Johnson Space Center, publiek domein; een cel van \emph{Emiliania huxleyi}, Alison R.~Taylor, University of North Carolina Wilmington Microscopy Facility, CC~BY~2.5; \emph{Orbulina universa}, H.~Takagi, K.~Kimoto, T.~Fujiki, H.~Saito, C.~Schmidt, M.~Kucera en K.~Moriya, CC~BY~4.0; diatomeeën uit de McMurdo Sound, Gordon T.~Taylor, Stony Brook University, NOAA Corps Collection, publiek domein (allemaal Wikimedia Commons).
  \item Hfst.~23 (diagenese en sedimentaire bekkens): kwartsoverwoekeringen in de Potsdam-zandsteen, M.~C.~Rygel, CC~BY-SA~4.0 (Wikimedia Commons). Zeespiegelstapel van R.~M.~Spratt en L.~E.~Lisiecki (2016), NOAA NCEI World Data Service for Paleoclimatology; compactieparameters van J.~G.~Sclater en P.~A.~F.~Christie (1980), zoals opgenomen in pyBacktrack (EarthByte).
  \item Hfst.~24 (de stratigrafische schaal): gouden spijker van het GSSP van het Selandien bij Zumaia, Josugoni, CC~BY-SA~4.0 (Wikimedia Commons); de grens Krijt--Paleogeen in het Trinidad Lake State Park, Nationalparks, CC~BY-SA~2.5 (Wikimedia Commons). Basisleeftijden en kleuren van de International Chronostratigraphic Chart, International Commission on Stratigraphy, CC~BY~4.0.
  \item Hfst.~25 (radiometrische dateringsmethoden): een zirkoonkorrel in kathodoluminescentie, Emmanuel Roquette, CC~BY-SA~4.0; een thermische-ionisatiemassaspectrometer, Radiogenic, CC0 (beide Wikimedia Commons). De ijkkromme IntCal20 van P.~J.~Reimer en anderen (2020), \emph{Radiocarbon} 62, CC~BY~4.0, zoals verspreid met het rintcal-pakket (M.~Blaauw).
  \item Hfst.~26 (structuur van de atmosfeer en de oceaan): het ophalen van een CTD en rosette in de Beringstraat, NASA Goddard Space Flight Center, publiek domein (Wikimedia Commons). Temperatuur en zoutgehalte uit de World Ocean Atlas 2023 (NOAA National Centers for Environmental Information).
  \item Hfst.~27 (straling en warmtetransport): de door de zon verlichte aarde van het Deep Space Climate Observatory, NASA Goddard Space Flight Center, publiek domein (Wikimedia Commons). Straling, oppervlakteflux en neerslag van de NCEP-NCAR Reanalysis~1 (E.~Kalnay en anderen, 1996), gegevens verstrekt door de NOAA PSL, Boulder, Colorado, VS, van hun website op https://psl.noaa.gov.
  \item Hfst.~28 (de algemene circulatie van de atmosfeer): een extratropische cycloon in de Golf van Alaska, NASA MODIS/LANCE, publiek domein; de intertropische convergentiezone uit GOES-11-wolkendata, NASA Earth Observatory, publiek domein (beide Wikimedia Commons). Wind en zeeniveaudruk van de NCEP-NCAR Reanalysis~1 (E.~Kalnay en anderen, 1996), gegevens verstrekt door de NOAA PSL, Boulder, Colorado, VS, van hun website op https://psl.noaa.gov.
  \item Hfst.~29 (de circulatie van de oceaan): de Golfstroom gezien via de zeeoppervlaktetemperatuur, NOAA NESDIS, publiek domein; een profilerende diepzeedrijver die wordt uitgezet, onbekende NOAA-fotograaf, publiek domein (beide Wikimedia Commons).
\end{itemize}
